%! Conditional Probability, Independence, and Simulation — Unit 3, Lesson 3.4 — Lesson Plan
% GRADUAL-RELEASE plan, timing 5 / 45 / 5 / 5: warm-up / guided notes and practice / spoken debrief
% / close and ASSIGN the homework. I do / we do in rows 1-4 (problems 1-10, ~25 min) -> we do the
% Guided Practice row (problems 11-14, ~20 min). Teacher-facing.
\documentclass[10pt]{article}
\usepackage{saar-article}
\usepackage{saar-boxes}
\usepackage{graphicx}   % -article does NOT load graphicx
\graphicspath{{images/}}
\newcommand{\TallMath}[1]{$\displaystyle #1\rule[-1.4em]{0pt}{3.2em}$}
\newcommand{\UnitNumberName}{Unit 3: Categorical Data and Probability \quad}
\newcommand{\LessonNumberName}{Lesson 3.4: Conditional Probability, Independence, and Simulation}

\begin{document}
\begin{center}
    {\Large\bfseries \CourseName \\
    {\small\bfseries \UnitNumberName \LessonNumberName}}
\end{center}

\begin{tcolorbox}[colback=frost, colframe=cerulean, boxrule=0.9pt, arc=2mm,
    left=2mm, right=2mm, top=1.5mm, bottom=1.5mm]
    \textbf{Primary Objective:} Students will find conditional probabilities from a two-way table
    and by formula, test two events for independence by comparing $P(A \mid B)$ with $P(A)$, read a
    pre-run simulation, and justify a decision with the conditional probability that actually
    answers the question.
    \\[2pt]
    \textbf{Standards:} PS.P.1d \quad$\cdot$\quad AFDA.DA.3b, AFDA.DA.3d \quad$\cdot$\quad
    AFDA.DA.3c (simulation as a probability model only) \quad$\cdot$\quad carries forward
    Lesson~3.3's multiplication rule \quad$\cdot$\quad folds in the former Lessons 3.5 and 3.6
    \\[2pt]
    \textbf{Lesson model:} \emph{Gradual release --- I do, we do, you do.} One Main Ideas / Notes
    table: four instruction rows (problems 1--10; you work 1, 3, 5, and 7, the class works the
    rest), then the Guided Practice row (11--14) with students holding the pen. \textbf{The
    homework is the individual practice}, assigned at Close \& Assign and done outside class. The
    debrief is a spoken cold check and one board item; there is no exit ticket.
    \\[2pt]
    \textbf{Conventions on the page:} the simulations are real runs shown as pre-made figures; no
    student runs one. Probabilities are decimals, rounded to two places in row~4 and the homework.
\end{tcolorbox}

\renewcommand{\arraystretch}{1.4}
\boxguard
\begin{skillbox}[Learning Targets \& Key Understandings]{goldbox}
\begin{tabularx}{\linewidth}{@{}X|X@{}}
\textbf{Learning targets (student language)} & \textbf{Key understandings (the ``why'')} \\[2pt]
\hline
\vspace{-6pt}
\begin{itemize}[leftmargin=1.1em, nosep, after=\vspace{2pt}]
  \item I can find a \textbf{conditional probability} from a two-way table by restricting to the
        ``given'' row or column. \hfill \textit{PS.P.1d}
  \item I can use $P(B \mid A) = P(A \text{ and } B)/P(A)$ on a Venn diagram, and I know it is $0$
        for mutually exclusive events. \hfill \textit{AFDA.DA.3b}
  \item I can test whether two events are \textbf{independent} by comparing $P(A \mid B)$ with
        $P(A)$. \hfill \textit{PS.P.1d}
  \item I can read a \textbf{simulation}, explain the \textbf{Law of Large Numbers}, and justify a
        decision. \hfill \textit{AFDA.DA.3c, 3d}
\end{itemize}
&
\vspace{-6pt}
\begin{itemize}[leftmargin=1.1em, nosep, after=\vspace{2pt}]
  \item \textbf{Target misconception:} $P(A \mid B) = P(B \mid A)$, swapping the condition ---
        ``most who passed used the app, so an app user will probably pass.'' \textbf{The rule that
        replaces it:} name the group the question is about, and put that group after the bar.
  \item \textbf{Second misconception:} dividing by the grand total, or by the wrong row or
        column's total. ``Given'' changes the whole.
  \item Dependent does not mean likely. In the simulation the app and passing are associated
        ($0.41$ for app users, about $0.16$ without), and an app user still probably does not pass.
  \item A short run can say almost anything. The Law of Large Numbers is why a long simulation
        can be used as a probability model.
\end{itemize}
\end{tabularx}
\end{skillbox}

\boxguard[24]
\begin{skillbox}[{Vocabulary, Concepts, \& Theorems --- taught directly in the guided notes}]{greenbox}
\small
\renewcommand{\arraystretch}{1.35}
\begin{tabularx}{\linewidth}{@{} >{\bfseries}p{3.4cm} X @{}}
Conditional probability & The probability of an event given that another happened, $P(A \mid B)$.
       ``Given $B$'' restricts to the $B$ group. \TallMath{P(B \mid A) = \frac{P(A \text{ and } B)}{P(A)}} \\
Independent (the test) & Knowing $B$ does not change the probability of $A$: $P(A \mid B) = P(A)$.
       If the two differ, the events are dependent (associated). \\
Simulation & Using a random model, such as a spreadsheet or calculator, to imitate a process many
       times and record the results. \\
Law of Large Numbers & As a process repeats more and more, the relative frequency of an event
       gets closer to its actual probability. \\
Mutually exclusive & Cannot both happen, so $P(A \text{ and } B) = 0$ and $P(B \mid A) = 0$. One
       printed line in row~2, nothing more (Lesson~3.3 taught the term). \\
\end{tabularx}
\end{skillbox}

\begin{fixedskillbox}[Lesson at a Glance --- \MeetingLength]{frost}
\renewcommand{\arraystretch}{1.25}
\begin{tabularx}{\linewidth}{@{}l c X X@{}}
\textbf{Phase} & \textbf{Min} & \textbf{Students} & \textbf{Teacher} \\
\hline
Warm-Up & 5 & Read the coding-course table, two percents of a part, compare decimals --- silent &
  Debrief item~2 aloud; write ``$84$ of $120$'' and ``$84$ of $100$'' on the board and leave them
  up all period \\
Guided Notes \& Practice & 45 & Fill the vocabulary box as each term is named; work rows 1--4
  (problems 1--10) with the pen in their hand; then the Guided Practice row, \emph{Spam Filter}
  (11--14), together & \textbf{I do / we do, row by row}: read the sentences, mark up the display,
  work problems 1, 3, 5, 7; the class works the rest (25) $\to$ \textbf{we do} Guided Practice
  11--14 (20) \\
Debrief & 5 & Pens down --- answer aloud, nothing to fill in & Put problem~14 back on the board
  ($0.90$ vs.\ $0.45$); ask the cold check \\
Close \& Assign & 5 & Write down the homework, the due date, and item~6 as the first item &
  Announce the homework and its form (DeltaMath + packet items 4, 5, 6), preview Lesson~3.5 \\
\end{tabularx}
\end{fixedskillbox}

\begin{fixedskillbox}[Warm-Up --- Activate Prior Knowledge (5 min)]{frost}
\begin{tabularx}{\linewidth}{@{}X X@{}}
\begin{minipage}[t]{\linewidth}\vspace{0pt}
    \textbf{The three items and what each seeds}
    \begin{itemize}[leftmargin=1.1em, itemsep=1pt, topsep=2pt]
      \item \textbf{1.} Read the table: $84$, $120$, $100$. \textbf{Seeds:} row~1 reopens this
            exact table.
      \item \textbf{2.} $84$ of $120$ is $70\%$; $84$ of $100$ is $84\%$. \textbf{Seeds:} problems
            1 and~2 --- the same part over two different wholes, which is the swap in miniature.
      \item \textbf{3.} $0.70 > 0.50$; $0.60 = 0.6$; $0.41 < 0.8$. \textbf{Seeds:} the comparison in
            row~3's test (problems 5 and~6) and row~4's $0.41$ against $0.80$.
    \end{itemize}
\end{minipage}
&
\begin{minipage}[t]{\linewidth}\vspace{0pt}
    \textbf{Running it}
    \begin{itemize}[leftmargin=1.1em, itemsep=1pt, topsep=2pt]
      \item \textbf{Debrief item~2 aloud.} Ask: same $84$, so why two percents? (Different wholes.)
            Leave both fractions on the board; row~1 points back to them.
      \item Item~3's last comparison is the tell: a student who says $0.41 > 0.8$ is comparing
            digit counts, and will misread row~4.
      \item Item~1 is reading only. Do not let it expand.
    \end{itemize}
\end{minipage}
\end{tabularx}
\end{fixedskillbox}

\boxguard
\begin{skillbox}[Hook --- before anyone sits down]{frost}
The hook is on \textbf{page 1 of the notes}, under the vocabulary box, and on the board: \textit{``Most
students who passed the course used our app, so a student who uses the app will probably pass.''}
\textbf{Ask: do you buy it?} Students circle \emph{Agree}, \emph{Disagree}, or \emph{Can't tell yet}
and write a reason \emph{before} anyone speaks; then take the hand vote and write the split on the
board. Most rooms agree. \textbf{Do not settle it.} Rows 1--3 use review sessions and class sections
on purpose, so no one computes the poster's numbers early. Say only that the class votes again at
problem~10, in the \emph{Simulate and Decide} row.
\end{skillbox}

\boxguard
\begin{skillbox}[Guided Notes \& Practice --- I do, then we do (45 min)]{frost}
\textbf{This is the whole lesson.} One Main Ideas / Notes table: four instruction rows (problems
1--10), then the Guided Practice row (11--14). \textbf{The rhythm in every row is the same}: read
the printed sentences aloud, mark up the display, work the first problem of the grid yourself, then
hand the rest to the room. The vocabulary box on page~1 is filled \emph{as each term is named}.
\begin{multicols}{2}
\textbf{I do / we do, row by row (about 25 min).}

\emph{Row 1, Two-Way Tables (problems 1--2).} The warm-up's table. Name \emph{conditional
probability}, fill the two ``whole'' boxes ($120$, $100$), and put your hand over everything outside
the App row. Work problem~1 ($84/120 = 0.70$). The class does 2 ($84/100 = 0.84$); ask which total
they used. Five minutes.

\emph{Row 2, The Formula (problems 3--4).} The same students as probabilities. Work problem~3 on
the Venn diagram ($0.42/0.60 = 0.70$) and point out that it matches problem~1. The class does 4
($0.08/0.40 = 0.20$). Read the mutually exclusive line and ask aloud for
$P(\text{Afternoon} \mid \text{Morning})$; it gets no more time than that. Five minutes.

\emph{Row 3, The Test (problems 5--6).} Name \emph{independent} and read the test. Work problem~5
($0.70$ against $0.50$: dependent). The class does 6 ($0.60$ and $0.60$: independent), the first
time ``equal'' has decided anything. Six minutes.

\emph{Row 4, Simulate and Decide (problems 7--10) --- the crux.} Name \emph{simulation}; tell the
class the figures are a finished spreadsheet run. Work problem~7: trace the line from $0.2$ after
$5$ users, up to about $0.57$ after $30$, then flat near $0.41$, and name the \emph{Law of Large
Numbers}. The class does 8 ($250/607 \approx 0.41$). \textbf{Problem~9 is the trap:} the same $250$
over $314$ gives $0.80$, so ask whether it is the same answer. Read the boxed sentence, then
\textbf{problem~10 is the crux: re-take the hook vote before anyone writes.} Dana's question names
app users, so $P(\text{Pass} \mid \text{App}) \approx 0.41$ answers it; the poster quoted the $0.80$.
Nine minutes.

\columnbreak
\textbf{We do (about 20 min).} The Guided Practice row, \emph{Spam Filter} (problems 11--14), a fresh
context with a table and no simulation. \textbf{Students hold the pen; you hold the questions.}
Problem~11 is $90/100 = 0.90$ and problem~12 is $90/200 = 0.45$: make students say each in words.
Problem~13 is the test ($0.90$ against $0.20$: dependent, so the filter is doing something).
\textbf{Problem~14 is the crux transferred}, and it comes last: take a hand count on Theo before
anyone answers. His question starts from a flagged email, so $0.45$ answers it, and $55\%$ of
flagged emails are real mail.

Circulate during Guided Practice and demand the reason, not the number:
\begin{itemize}[leftmargin=1.1em, nosep]
  \item \textbf{Problems 11--12:} ``Point to the group you restricted to. What is the whole?''
  \item \textbf{Problem 13:} ``What two numbers did you compare? What would equal have meant?''
  \item \textbf{Problem 14:} ``What does Theo know about the email in front of him? So what goes
        after the bar?'' A student who answers $0.90$ has not left the hook vote.
\end{itemize}
While circulating, pick one paper to put up in the debrief: a problem~14 that named
$P(\text{Spam} \mid \text{Flagged})$. \textbf{This circulation is your formative read.}
\end{multicols}
\end{skillbox}

\boxguard[24]
\begin{skillbox}[Debrief --- whole class, spoken (5 min)]{frost}
Pens down, papers up. \textbf{This debrief is spoken} --- it lives on the board and in the room.
\begin{multicols}{2}
\begin{enumerate}[leftmargin=1.3em, nosep, itemsep=3pt]
  \item \textbf{The one board item: problem~14.} Put up the paper you picked, $0.90$ beside $0.45$,
        and make a student say the sentence: \emph{the question starts from a flagged email, so
        flagged goes after the bar.} Point back at the warm-up's $84$ of $120$ and $84$ of $100$.
\end{enumerate}
\columnbreak
\textbf{Check for understanding --- ask it cold, whole class.} \emph{``Most students on the honor
roll eat breakfast. Jay says a student who eats breakfast will probably make the honor roll. What
did he do?''} Thirty seconds of think time, then cold-call three students. Correct answer: he
\textbf{swapped the condition}. He knows $P(\text{Breakfast} \mid \text{Honor roll})$ is high; his
claim needs $P(\text{Honor roll} \mid \text{Breakfast})$, which he has not found. \emph{This replaces
the exit ticket.}

\textbf{Formative read.} Sort the answers into three piles: \textbf{(a)} named both conditionals in
``given'' language; \textbf{(b)} said he has it backwards but could not say which group goes after
the bar; \textbf{(c)} agreed with Jay. \textbf{If pile (c) runs past a third of the room, open
Lesson~3.5 by re-running problems 9--10.}
\end{multicols}
\end{skillbox}

\boxguard
\begin{skillbox}[Homework --- scored 1 to 5, assigned today, due after two study halls]{goldbox}
Six items on a \textbf{free-throw simulation and a coach's decision} ($1{,}000$ simulated late-game
trips; Riley takes $800$, Sam $200$). \textbf{1} and \textbf{2} are the \textbf{contrast pair}: the
same $286$ over two wholes, $P(\text{Made both} \mid \text{Riley}) \approx 0.36$ and
$P(\text{Riley} \mid \text{Made both}) \approx 0.69$; \textbf{3} is \emph{the spiral item}, Lesson~3.3's
multiplication rule ($0.64$ and $0.36$); \textbf{4} is the justification, the independence test
($0.64$ against about $0.41$); \textbf{5} is the interpretation, the Law of Large Numbers ($0.80$
after $20$ trips, $0.64$ after $200$); \textbf{6} is \emph{the crux head-on}, the assistant's swap and
the decision (Sam should shoot).

\textbf{Assigned at Close \& Assign and done outside class.} \textbf{Scoring:} 12 points, which the
grading report turns into the score of 1 to 5. Items 4--6 are rubric boxes (see \texttt{rubrics/});
items 1--3 are unboxed, a worked reference for the DeltaMath set.

\textbf{Packet, Desmos, or DeltaMath --- decided here: DeltaMath + justifications.} Items 1--2 are
conditional probabilities read from a table and item~3 is the multiplication rule for independent
shots, which DeltaMath drills and checks; items 4, 5, and~6 are the independence comparison and what
it means, the Law of Large Numbers argument, and the assistant's reversed conditional --- written
reasons it cannot grade. So assign \textbf{DeltaMath} --- Conditional Probability from a Table
$\cdot$ Conditional Probability from Table (L2) $\cdot$ Probability of Independent \& Dependent
Events $\cdot$ Determine Independence from a Table $\cdot$ Probability from Simulation --- for the
computation (items 1--3 stay in the packet as a worked reference), and students \textbf{turn in packet
items 4, 5, and~6}. Still 12 points: 2 per turned-in item, the DeltaMath score scaled to the other
6. Say so aloud at Close \& Assign.

\textbf{Preview of the next class.} Lesson~3.5, Counting and Probability with Counting.

\textbf{Connections \& Big Ideas.} \textit{Unit 3 core idea:} organize categorical information, then
use probability rules to make decisions. \textit{Standards covered:} PS.P.1d (conditional
probabilities to determine association or independence), AFDA.DA.3b (conditional probability,
including mutually exclusive events), AFDA.DA.3c (simulation as a probability model), AFDA.DA.3d
(interpret simulation results to make and justify decisions). \textit{Reaches back to:} Lesson~3.1's
two-way tables and Lesson~3.3's addition and multiplication rules. \textit{Spirals forward to:}
Lesson~3.5's counting, then the Unit~3 review and test.
\end{skillbox}

\boxguard
\begin{skillbox}[Watch For (while circulating)]{redbox}
\begin{itemize}[leftmargin=1.2em, nosep]
  \item \textbf{Swapping the condition} (the hook, problems 9, 10, 12, 14; HW~2, 6 --- the target
        misconception). Probe: \emph{``Who is the question about? Put that group after the bar.''}
  \item \textbf{Dividing by the grand total} (problems 1, 2, 8; HW~1). Probe: \emph{``Cover
        everyone who is not in the given group. How many are left?''}
  \item \textbf{Comparing the wrong pair in the test} (problems 5, 6, 13; HW~4): comparing
        $P(A \mid B)$ with $P(B \mid A)$. Probe: \emph{``The test asks whether knowing changes it.
        What is the chance before you know anything?''}
  \item \textbf{Dependent read as ``probably''} (problem~10). Probe: \emph{``The app raises the
        chance. Raises it to what?''}
  \item \textbf{Trusting the early part of the line} (problem~7; HW~5). Probe: \emph{``How many
        users is that based on? What did the line do next?''}
\end{itemize}
Cold-call prompts: \emph{``Give me two events at this school where the two conditionals are very
different.''} \quad \emph{``What would the table look like if App and Morning were not independent?''}
\end{skillbox}

\boxguard
\begin{skillbox}[Close \& Assign (5 min)]{goldbox}
\textbf{The homework is assigned here and done outside class --- it is not started in class.} Hand
the launch: six items on the free-throw simulation, scored 1 to 5, due at the start of the first
class after two study halls. Say the form: \textbf{DeltaMath} for the conditional-probability,
independence, and simulation practice, and \textbf{packet items 4, 5, and~6 turned in}. \textbf{Tell students to do item~6 first if they are
short on time}, and that items 1 and~2 are a pair: read both before answering either. Then close on
what changed: \emph{``You walked in able to read a table. You walk out knowing that `given' changes
the whole, and that the first thing to ask about any probability claim is: given what?''}
\textbf{Preview:} Lesson~3.5, Counting and Probability with Counting.
\end{skillbox}

% Teacher-only prose lives HERE, never in a _key.
\begin{teachernote}[Warm-Up]
Five minutes, silent, timed. The table is row~1's, so row~1 can move fast. \textbf{Item~2 is the one
to read while collecting}: the same $84$ gives $70\%$ and $84\%$, and a student who can say why has
the idea before it has a name. Leave both fractions on the board; the debrief points back at them.
Note any student who wrote $0.41 > 0.8$ and check on them at problems 8 and~9.
\end{teachernote}

\begin{teachernote}[Guided Notes \& Practice]
\textbf{Pace it 25 / 20, then 5 for the debrief and 5 to assign the homework.} You work 1, 3, 5, and
7. Rows 1--3 must move, because \textbf{the time belongs in row~4}. Rows 1--3 deliberately avoid
passing, so that the poster's two numbers appear for the first time side by side at problems 8
and~9. \textbf{Re-take the vote at problem~10 before anyone writes}; if students compute first, the
vote is empty. Keep the verdict honest: the app and passing \emph{are} associated ($64/393 \approx
0.16$ without the app, against $0.41$ with it), so the poster is wrong about ``probably,'' not about
the app being worth anything. The notes do not claim the app causes passing, and a sharp student may
ask; say that a simulation built from past rates cannot settle cause.

\textbf{About the figures.} Both are one real run: a spreadsheet-style model of $1{,}000$ students with
$P(\text{App}) = 0.60$, $P(\text{Pass} \mid \text{App}) = 0.40$, $P(\text{Pass} \mid \text{No app}) =
0.15$. Those inputs are not printed, so the students read the output rather than the recipe. Students
do not run a simulation in this lesson. If time is short in Guided Practice, protect 12 and~14;
problem~13 can be said aloud.
\end{teachernote}

\begin{teachernote}[Homework]
\textbf{Scored 1 to 5 from 12 points, due at the start of the first class after two
study halls; assigned at the close, not started in class.} Items 4--6 follow their rubrics.
\textbf{Item~6 is the one to read
first.} Sort it into three piles: named the swap and compared $0.36$ with $0.64$ (ready); chose Sam
without naming the swap (needs the language); chose Riley (the misconception survived). If more than
a third land in the last pile, open Lesson~3.5 with problems 9--10. The simulation is a real run: Riley
$0.6$ a shot for $800$ trips, Sam $0.8$ a shot for $200$, two shots a trip.

\textbf{DeltaMath + justifications is the call for 3.4.} DeltaMath (exact topic names): Conditional
Probability from a Table $\cdot$ Conditional Probability from Table (L2) $\cdot$ Probability of
Independent \& Dependent Events $\cdot$ Determine Independence from a Table $\cdot$ Probability from
Simulation. Students turn in packet items 4, 5, and~6. Score the three turned-in items by their
rubrics (item~6 first, then 4, then 5); scale the DeltaMath score to the other 6 points. The items
1--2 contrast pair (same numerator, different condition) moves to the two conditional-from-a-table
topics: a student who gets the same value for both directions there has not restricted. Item~6 still
requires both conditionals in writing, so the swap is still checked on paper.
Collect the turned-in items at the start of the first class after two study halls.
\end{teachernote}
\end{document}
